\clearpage
\section{Αναλυτικό mapping της εκτέλεσης}
\label{app:ch2-execution-map}

Στο κύριο κείμενο δόθηκε η γενική εικόνα. Εδώ καταγράφονται πιο αναλυτικά οι
βασικές εκτελεστικές διαδρομές με αναφορά στα modules που συμμετέχουν.

\subsection*{Single run για τροχιά και phase portrait}
\begin{itemize}
\item Η διεπαφή διαβάζει τα widgets και συγκροτεί τα objects
\texttt{SystemConfig}, \texttt{IntegrationConfig}, \texttt{InitialConditions}
και \texttt{SolverTolerances}.
\item Η \texttt{solve\_cached} στο \path{app/cache.py} αποφασίζει ποια
αριθμητική μέθοδος θα χρησιμοποιηθεί.
\item Για adaptive runs καλείται το \path{core/solver.py}. Για
symplectic runs καλείται το \path{core/symplectic_solver.py}.
\item Το αποτέλεσμα επιστρέφει ως πίνακες \(t\) και \(y\), οι οποίοι
χρησιμοποιούνται από το plotting layer.
\end{itemize}

\subsection*{Μεμονωμένος υπολογισμός Lyapunov}
\begin{itemize}
\item Η διεπαφή ορίζει \(t_{\mathrm{tr}}\), \(t_{\mathrm{m}}\),
\(\Delta t\) και QR interval.
\item Η \texttt{compute\_lyapunov\_cached} προετοιμάζει RHS, Jacobian και
αριθμητικές επιλογές.
\item Ο πυρήνας \texttt{core/lyapunov.py} ολοκληρώνει το διευρυμένο σύστημα,
κάνει περιοδικό QR και επιστρέφει το φάσμα.
\end{itemize}

\subsection*{Bifurcation και Lyapunov sweeps}
\begin{itemize}
\item Το \path{app/ui/bifurcation_tab.py} συγκεντρώνει τις ρυθμίσεις sweep
και συντηρεί το state για continue workflows.
\item Για bifurcation sweep η βασική είσοδος περνά στο \path{app/sweep.py},
το οποίο τελικά χρησιμοποιεί το \path{core/poincare_sweep.py}.
\item Για Lyapunov sweep η ροή περνά από το
\path{app/logic/lyapunov_sweep.py} και καταλήγει στο
\path{core/lyapunov.py}.
\end{itemize}

Τα παραπάνω αντιστοιχούν στις βασικές πηγές κώδικα που αναφέρθηκαν στο
Κεφάλαιο~\ref{ch:computational-implementation}
\cite{nlds_src_app_main,nlds_src_ui_bifurcation_tab,nlds_src_app_cache,nlds_src_app_sweep,nlds_src_logic_lyapunov_cached,nlds_src_logic_lyapunov_sweep,nlds_src_core_solver,nlds_src_core_poincare,nlds_src_core_symplectic,nlds_src_core_lyapunov}.

\section{Ανάλυση πολυπλοκότητας}
\label{app:ch2-complexity}

\subsection{Ορισμοί συμβόλων}

\begin{table}[H]
\centering
\caption{Συμβολισμός που χρησιμοποιείται στην ανάλυση πολυπλοκότητας}
\label{tab:ch2-complexity-symbols}
\renewcommand{\arraystretch}{1.15}
\begin{tabular}{|l|p{0.64\textwidth}|}
\hline
\textbf{Σύμβολο} & \textbf{Ερμηνεία} \\
\hline
\(n\) & Διάσταση του διανύσματος κατάστασης. \\
\hline
\(N\) & Πλήθος fixed-step βημάτων, περίπου \((t_f-t_0)/\Delta t\). \\
\hline
\(P\) & Πλήθος τιμών παραμέτρου σε sweep. \\
\hline
\(H\) & Μέγιστο πλήθος διατηρούμενων hits ανά παράμετρο. \\
\hline
\(N_{\mathrm{QR}}\) & Πλήθος QR ορθοκανονικοποιήσεων στον υπολογισμό Lyapunov. \\
\hline
\(C_f(n)\) & Κόστος μίας αξιολόγησης του RHS. \\
\hline
\(C_J(n)\) & Κόστος μίας αξιολόγησης του Jacobian. \\
\hline
\end{tabular}
\end{table}

\subsection{Μοναδική τροχιά}

Για fixed-step RK4 ο χρόνος κλιμακώνεται γραμμικά με το πλήθος βημάτων:
\begin{equation}
T_{\mathrm{RK4}} = O\!\left(N\,C_f(n)\right).
\end{equation}
Η πλήρης αποθήκευση της τροχιάς έχει τάξη
\begin{equation}
M_{\mathrm{traj}} = O(nN).
\end{equation}

Για adaptive ολοκληρωτές τύπου \texttt{solve\_ivp}, το πραγματικό κόστος
εξαρτάται από το πλήθος εσωτερικών evaluations που επιβάλλονται από τις
ανοχές:
\begin{equation}
T_{\mathrm{IVP}} = O\!\left(N_f\,C_f(n)\right).
\end{equation}

\subsection{Bifurcation sweep}

Στην απλή fixed-step προσέγγιση, ένα sweep παραμέτρου έχει προσεγγιστικό κόστος
\begin{equation}
T_{\mathrm{sweep}} = O(PN).
\end{equation}
Αν χρησιμοποιούνται events και early stop, ο πραγματικός χρόνος μπορεί να είναι
μικρότερος, επειδή η συλλογή σταματά μόλις συμπληρωθεί ο αναγκαίος αριθμός
σημείων.

Η μνήμη που χρειάζεται για τα αποθηκευμένα sweep points είναι της τάξης
\begin{equation}
M_{\mathrm{sweep}} = O(PH).
\end{equation}

\subsection{Lyapunov φάσμα}

Με analytic Jacobian, το κόστος ενός run για πλήρες φάσμα γράφεται
προσεγγιστικά ως
\begin{equation}
T_{\mathrm{lya}}
=
O\!\left(N\,[C_f(n)+C_J(n)+n^3]\right)
+ O\!\left(N_{\mathrm{QR}}\,n^3\right).
\end{equation}

Με finite-difference Jacobian, το κόστος αυξάνεται επειδή κάθε στήλη του
Jacobian προσεγγίζεται από επιπλέον αξιολογήσεις του RHS:
\begin{equation}
T_{\mathrm{lya,FD}}
=
O\!\left(N\,[n\,C_f(n)+n^3]\right).
\end{equation}

\subsection{Παράλληλα sweeps}

Όταν τα runs για διαφορετικές τιμές παραμέτρου είναι ανεξάρτητα, ο ιδανικός
wall-clock χρόνος γράφεται σχηματικά ως
\begin{equation}
T_{\mathrm{wall}} \approx
O\!\left(\frac{P}{W}\cdot \mathrm{cost\_per\_run}\right)
+ O(T_{\mathrm{overhead}}),
\end{equation}
όπου \(W\) είναι ο αριθμός workers. Στο continuation mode τέτοια
παραλληλοποίηση δεν είναι γενικά δυνατή, επειδή κάθε run εξαρτάται από το
τελικό state του προηγούμενου.

\section{Διαχείριση μνήμης σε μεγάλα runs}
\label{app:ch2-memory}

\subsection{Bounded trajectory storage}

Η τρέχουσα υλοποίηση δεν αποθηκεύει υποχρεωτικά όλα τα σημεία μιας μεγάλης
τροχιάς. Αν το θεωρητικό πλήθος δειγμάτων είναι
\begin{equation}
N_{\mathrm{nom}}=
\left\lfloor\frac{t_f-t_0}{\Delta t}\right\rfloor + 1,
\end{equation}
τότε με ενεργό \texttt{max\_store\_steps} αποθηκεύονται τελικά
\begin{equation}
N_{\mathrm{store}}=
\min\!\left(N_{\mathrm{nom}},N_{\mathrm{store,max}}\right).
\end{equation}

Άρα το raw αποτύπωμα μνήμης των arrays εξόδου μπορεί να εκτιμηθεί από
\begin{equation}
M_{\mathrm{traj,out}}
\approx
8\,(n+1)\,N_{\mathrm{store}}
\quad \text{bytes},
\end{equation}
καθώς αποθηκεύονται ένας πίνακας χρόνου και \(n\) πίνακες κατάστασης σε
\texttt{float64}.

\subsection{Decimation για plotting}

Για τα γραφήματα χρησιμοποιείται επιπλέον downsampling. Αν το plot δείχνει
το πολύ \(N_{\mathrm{plot,max}}\) σημεία, τότε η μνήμη και ο χρόνος render
παραμένουν ελεγχόμενα με
\begin{equation}
N_{\mathrm{plot}}=
\min\!\left(N_{\mathrm{store}},N_{\mathrm{plot,max}}\right).
\end{equation}

\subsection{Bounded-memory bifurcation sweeps}

Στο interactive bifurcation plot διατηρούνται:
\begin{itemize}
\item ένα πρόσφατο παράθυρο σημείων σε πλήρη ανάλυση,
\item και ένα reservoir sample από παλαιότερα σημεία.
\end{itemize}

Αν \(N_{\mathrm{recent}}\) είναι το πλήθος των πρόσφατων σημείων και
\(N_{\mathrm{res}}\) το μέγεθος του reservoir, τότε το αποτύπωμα των
ζευγών \((x,y)\) είναι προσεγγιστικά
\begin{equation}
M_{\mathrm{bif,plot}}
\sim
16\,(N_{\mathrm{recent}}+N_{\mathrm{res}})
\quad \text{bytes},
\end{equation}
αγνοώντας το μικρό σταθερό overhead των δομών Python.

\section{Numba/JIT και compiled πυρήνες}
\label{app:ch2-numba}

\begin{toolbox}{Interpreted και compiled path στην εφαρμογή}{lam2015_numba,numba_jit_docs,numba_architecture_docs}
\textbf{Interpreted path:} η ορχήστρωση, η διεπαφή και τα fallback numerical
μονοπάτια τρέχουν κανονικά στον Python interpreter.

\textbf{Compiled path:} επιλεγμένα numerical kernels μεταγλωττίζονται δυναμικά
με Numba σε machine code και εκτελούνται εκτός του interpreter.
\end{toolbox}

Στην πράξη, το project χρησιμοποιεί compiled πυρήνες κυρίως για:
\begin{itemize}
\item fixed-step RK4 ολοκλήρωση,
\item symplectic Forest--Ruth ολοκλήρωση,
\item υπολογισμό φάσματος Lyapunov,
\item και ορισμένες διαδρομές Poincar\'e sweep.
\end{itemize}

Οι υλοποιήσεις αυτές βρίσκονται κυρίως στα
\path{core/numba_backend/} και \path{app/numba_custom.py}
\cite{nlds_src_core_numba,nlds_src_app_numba_custom}.

\begin{figure}[H]
\centering
\includegraphics[width=0.94\textwidth]{figures/ch2/numba_jit_llvm_pipeline.png}
\caption{Σχηματική ροή μεταγλώττισης μέσω Numba και LLVM προς native machine code
\cite{numba_architecture_docs,llvmlite_docs,llvm_langref}.}
\label{fig:ch2-numba-llvm-pipeline}
\end{figure}

\subsection{Πότε ενεργοποιείται}

Η ενεργοποίηση του compiled path εξαρτάται από δύο πράγματα:
\begin{itemize}
\item αν είναι διαθέσιμη η Numba στο περιβάλλον,
\item και αν το συγκεκριμένο numerical μονοπάτι υποστηρίζεται από τα
διαθέσιμα kernels.
\end{itemize}

Τυπικά, τα adaptive \texttt{RK45}/\texttt{DOP853} paths παραμένουν στην SciPy
διαδρομή, ενώ τα fixed-step runs είναι οι πιο φυσικοί υποψήφιοι για compiled
εκτέλεση.

\subsection{Fallback συμπεριφορά}

Το έργο δεν βασίζεται απόλυτα σε μία compiled διαδρομή. Αν ένα JIT path
αποτύχει, το orchestration layer επιστρέφει σε Python fallback χωρίς να αλλάζει
η διεπαφή του χρήστη. Αυτό κάνει την εφαρμογή πιο ανθεκτική σε custom μοντέλα
και σε διαφορετικά execution environments
\cite{nlds_src_app_cache,nlds_src_logic_lyapunov_cached,nlds_src_logic_lyapunov_sweep}.
